\documentclass[../../../include/open-logic-section]{subfiles}

\begin{document}

\olfileid{sth}{z}{sep}
\olsection{వేరుచేయడం}

అమాయక ధర్మసంగ్రహం స్థానంలో
ఉండే ఒక సూత్రంతో
మొదలుపెడదాం:

\begin{axiom}[వేరుచేయడం స్వీకృత పథకం] ప్రతి $\phi(x)$ సూత్రానికి
ఇది ఒక స్వీకృతం: ఏ $A$కైనా, $\Setabs{x \in A}{\phi(x)}$ సమితి ఉంటుంది.
\end{axiom}

ఇది ఒకే స్వీకృతం కాదని
గమనించండి. ఇది స్వీకృతాల
\emph{పథకం}. వేరుచేయడం
స్వీకృతాలు \emph{అనంతంగా}
ఉంటాయి; ప్రతి $\phi(x)$
సూత్రానికి ఒకటి. ఈ పథకాన్ని
సమానార్థకంగా ఈ కింది
విధంగా కూడా రాయవచ్చు;
సాధారణంగా అలాగే రాస్తారు:

\begin{defish}
``$S$'' లేని ప్రతి
$\phi(x)$ సూత్రానికి
ఇది ఒక స్వీకృతం:
\[
	\forall A \exists S \forall x(x \in S \liff (\phi(x) \land x \in A)).
\]
\end{defish}

\olref[sth][][]{part} మొదట
చెప్పిన సంప్రదాయం ప్రకారం,
వేరుచేయడం స్వీకృతాల్లోని
$\phi$ సూత్రాల్లో పరామితులు
ఉండవచ్చు.\footnote{దీని అర్థం
వివరణకు \olref[sfr][infinite][induction]{natinductionschema}
వెంటనే తరువాతి చర్చ చూడండి.}

మనం చెప్పుకుంటున్న సమితుల
సంచిత-దశలవారీ భావన నుంచి
వేరుచేయడాన్ని వెంటనే
సమర్థించవచ్చు. ఎందుకో
చూడడానికి $A$ ఒక సమితి
అనుకుందాం. అప్పుడు
\stageshier\ ప్రకారం, $A$
ఏదో ఒక దశ~$S$లో
ఏర్పడుతుంది. $A$ ఆ
దశ~$S$లో ఏర్పడింది కాబట్టి,
$A$ యొక్క మూలకాలన్నీ $S$
దశకు ముందే ఏర్పడ్డాయి
(\stagesacc\ ప్రకారం).
ముఖ్యంగా, $A$లో మూలకాలుగా
ఉండి $\phi$ను కూడా
సంతృప్తిపరిచే సమితులన్నింటినీ
పరిగణించండి. అవి కూడా
$S$ దశకు ముందే ఏర్పడ్డాయని
స్పష్టం. కాబట్టి ఆ సమితులు
$\Setabs{x \in A}{\phi(x)}$
అనే సమితిగా $S$ దశలోనే
ఏర్పడతాయి (\stagesacc\ ప్రకారం).

అమాయక ధర్మసంగ్రహానికి
భిన్నంగా, ఇది రసెల్
వైరుధ్యాన్ని తప్పిస్తుంది.
ఎందుకంటే $\Setabs{x}{x \notin
x}$ సమితి ఉంటుందని
నేరుగా ప్రకటించలేం.
దానికి బదులు, ఏదైనా
$A$ సమితి \emph{ఇచ్చినప్పుడు},
$R_A = \Setabs{x \in A}{x \notin x}$
సమితి ఉంటుందని చెప్పగలం.
కానీ దీని నుంచి వచ్చేది
$R_A \notin R_A$ మరియు
$R_A \notin A$ మాత్రమే;
ఇవి ఆందోళన కలిగించేవి కావు.

అయితే వేరుచేయడానికి ఒక
తక్షణ, ఆశ్చర్యకరమైన
పర్యవసానం ఉంది:

\begin{thm}\ollabel{thm:NoUniversalSet}
\emph{విశ్వవ్యాప్త} సమితి లేదు;
అంటే $\Setabs{x}{x = x}$ ఉండదు.
\end{thm}

\begin{proof}
వైరుధ్యాన్ని చూపడానికి,
$V$ విశ్వవ్యాప్త సమితి
అనుకుందాం. అప్పుడు
వేరుచేయడం ప్రకారం
$R =
\Setabs{x \in V}{x \notin x} = \Setabs{x}{x \notin x}$
ఉంటుంది; ఇది రసెల్
వైరుధ్యానికి దారితీస్తుంది.
\end{proof}

విశ్వవ్యాప్త సమితి లేకపోవడం---నిజానికి
సమితుల స్థాయిక్రమం
అంతంలేకుండా కొనసాగడం---సంచిత-దశలవారీ
భావనకు పునాదిగా ఉన్న
అత్యంత ముఖ్యమైన ఆలోచనల్లో
ఒకటి. దీనికి సహజమైన
మరో మార్గంలోనూ చేరుకోవచ్చని
చూడడం విలువైనది.
విశ్వవ్యాప్త సమితి తనకే
\tetoken{మూలకం}{!!a{element}}
అయి ఉండాలి. కానీ మన
సంచిత-దశలవారీ భావనలో,
ప్రతి సమితీ తన
\tetoken{మూలకాలన్నింటి}{!!{element}s}
తరువాత స్థాయిక్రమంలో
వచ్చే తొలి దశలో కనిపిస్తుంది.
అందువల్ల \emph{ఏ} సమితీ
తనకే మూలకం కాజాలదు.
తనకే మూలకమైన సమితి,
తాను మొదట కనిపించిన
దశకు తరువాత మాత్రమే
మొదట కనిపించాల్సి వస్తుంది;
ఇది అసంబద్ధం. (సమితి
``స్థాయి'' భావాన్ని ఉపయోగించి
ఈ వివరణను మరింత కచ్చితంగా
ఎలా చెప్పాలో
\olref[sth][spine][rank]{defnsetrank}లో
చూస్తాం. అయితే దానికి
ముందు మరికొన్ని స్వీకృతాలు
అవసరం.)

వేరుచేయడం, సమితుల సమానత్వ
సూత్రం కలిపి ఇచ్చే మరికొన్ని
పర్యవసానాలు ఇవి.

\begin{prop}\ollabel{prop:emptyexists}
ఏదైనా సమితి ఉంటే,
$\emptyset$ ఉంటుంది.
\end{prop}

\begin{proof}
$A$ ఒక సమితి అయితే,
వేరుచేయడం ద్వారా
$\emptyset = \Setabs{x \in A}{x \neq x}$
ఉంటుంది.
\end{proof}

\begin{prop}
ఏ సమితులు $A$, $B$కైనా
$A \setminus B$ ఉంటుంది.
\end{prop}

\begin{proof}
వేరుచేయడం ద్వారా
$A \setminus B = \Setabs{x \in A}{x \notin B}$
ఉంటుంది.
\end{proof}

(దాదాపు) ఏవైనా సమితుల
సర్వసామాన్య ఛేదనలు కూడా
ఉంటాయని తేలుతుంది:

\begin{prop}\ollabel{prop:intersectionsexist}
$A \neq \emptyset$ అయితే,
$\bigcap A = \Setabs{x}{(\forall y \in A)x \in y}$
ఉంటుంది.
\end{prop}

\begin{proof}
$A \neq \emptyset$ అనుకుందాం;
అప్పుడు ఏదో $c \in A$
ఉంటుంది. కాబట్టి
$\bigcap A =
\Setabs{x}{(\forall y \in A)x \in y} = \Setabs{x \in c}{(\forall y \in
A)x \in y}$; వేరుచేయడం
ప్రకారం ఇది ఉంటుంది.
\end{proof}

అయితే $A \neq \emptyset$
అనే నిబంధనను గమనించండి;
లేకపోతే $\bigcap
\emptyset$ శూన్యసత్యం వల్ల
విశ్వవ్యాప్త సమితి అవుతుంది;
అది \olref{thm:NoUniversalSet}కు
విరుద్ధం.

\end{document}
